\section{Introduction}

The lightest bound state in quantum chromodynamics (QCD), the pion, plays a fundamental role, since it is the Nambu-Goldstone boson of dynamical chiral symmetry breaking (DCSB).
Studies of pion and kaon structure reveal the physics of DCSB, help to reveal the relative impact of DCSB versus the chiral symmetry breaking by the quark masses, and are important to understand nonperturbative QCD.
Studying the pion parton distribution functions (PDFs) is important to characterize the structure of the pion and further understand DCSB and nonperturbative QCD.
Currently, our knowledge of the pion PDFs is less than the nucleon PDFs, because there are fewer experimental data sets, especially for the sea-quark and gluon distributions.
The future U.S.-based Electron-Ion Collider (EIC)~\cite{Accardi:2012qut}, planned to be built at Brookhaven National Lab, will further our knowledge of  pion structure~\cite{Arrington:2021biu,Aguilar:2019teb}.
%
In China, a similar machine, the Electron-Ion Collider in China (EicC)~\cite{Anderle:2021wcy}, is also planned to make impacts on the pion gluon and sea-quark distributions.
%
In Europe, the Drell-Yan and $J/\psi$-production experiments from COMPASS++/AMBER~\cite{Denisov:2018unj} will aim at improving our knowledge of both the pion gluon and quark PDFs.

Global analyses of pion PDFs mostly rely on Drell-Yan data.
The early studies of pion PDFs were based mostly on pion-induced Drell-Yan data and use $J/\psi$-production data or direct photon production to constrain the pion gluon PDF~\cite{Owens:1984zj,Aurenche:1989sx,Sutton:1991ay,Gluck:1991ey,Gluck:1999xe}.
There are more recent studies, such as the work by Bourrely and Soffer~\cite{Bourrely:2018yck}, that extract the pion PDF based on Drell-Yan $\pi^+W$ data.
JAM Collaboration~\cite{Barry:2018ort,Cao:2021aci} uses a Monte-Carlo approach to analyze the Drell-Yan $\pi A$ and leading-neutron electroproduction data from HERA to reach the lower-$x$ region, and revealed that gluons carry a significantly higher momentum fraction (about $30\%$) in the pion than had been inferred from Drell-Yan data alone.
The xFitter group~\cite{Novikov:2020snp} analyzed Drell-Yan $\pi A$ and photoproduction data using their open-source QCD fit framework for PDF extraction and found that these data can constrain the valence distribution well but are not sensitive enough for the sea and gluon distributions to be precisely determined.
The analysis done Ref.~\cite{Chang:2020rdy} suggests that the pion-induced $J/\psi$-production data has additional constraint on pion PDFs, particularly in the pion gluon PDF in the large-$x$ region.
All in all, the pion valence-quark distributions are better constrained
than the gluon distribution from the global analysis of experimental data.
While waiting for more experimental data sets, the study of the pion gluon distribution from theoretical side can provide useful information for the experiments.

The pion gluon PDF is rarely studied using continuum-QCD phenomenological models or through lattice-QCD (LQCD) simulations.
Most model studies only predict the pion valence-quark distribution~\cite{Nam:2012vm,Watanabe:2016lto,Watanabe:2017pvl,Hutauruk:2016sug,Lan:2019vui,Lan:2019rba,deTeramond:2018ecg,Watanabe:2019zny,Han:2018wsw,Chang:2014lva,Chang:2014gga,Chen:2016sno,Shi:2018mcb,Bednar:2018mtf,Ding:2019lwe},
but the gluon and sea PDFs are predicted by the Dyson-Schwinger equation (DSE) continuum approach~\cite{Freese:2021zne}. The prediction of the pion gluon PDF in DSE, based on an implementation of rainbow-ladder truncation of DSE, is consistent with the JAM pion gluon PDF result~\cite{Barry:2018ort,Cao:2021aci} within two sigma.
LQCD provides an first-principles calculations to improve our knowledge of nonperturbative pion gluon structure;
however, there have been only two efforts to determine the first moment of pion gluon PDF~\cite{Meyer:2007tm,Shanahan:2018pib}.
An early calculation in 2000 using quenched QCD predicted $\langle x \rangle_g=0.37(8)(12)$, using Wilson fermion action with a lattice spacing $a=0.093$~fm, lattice size $L^3\times T=24^4$, a large 890-MeV pion mass and 3,066 configurations at $\mu^2=4\text{ GeV}^2$~\cite{Meyer:2007tm}.
A more recent study in 2018 using $N_f=2+1$ clover fermion action  with a lattice spacing $a=0.1167(16)$~fm, larger lattice size $32^3\times 96$, 450-MeV pion mass, and 572,663 measurements, gave a larger first-moment result, $\langle x \rangle_g=0.61(9)$ at $\mu^2=4\text{ GeV}^2$~\cite{Shanahan:2018pib}.
%
In principle, a series of moments can be used to reconstruct the PDF.
Although there are calculations of the first moment of the pion gluon PDF, there is little chance that sufficient higher moments of the pion gluon PDF can be obtained to perform such a reconstruction.
A direct lattice calculation of the $x$-dependence of the pion gluon PDF is needed.

In recent years, there has been an increasing number of calculations of $x$-dependent hadron structure in lattice QCD, following the proposal of Large-Momentum Effective Theory (LaMET)~\cite{Ji:2013dva,Ji:2014gla,Ji:2017rah}.
The LaMET method calculates lattice quasi-distribution functions, defined in terms of matrix elements of equal-time and spatially separated operators, and then takes the infinite-momentum limit to extract the lightcone distribution.
The quasi-PDF can be related to the $P_z$-independent lightcone PDF through a factorization theorem that factors from it a perturbative matching coefficient with corrections suppressed by the hadron momentum~\cite{Ji:2014gla}.
The factorization can be calculated exactly in perturbation theory~\cite{Ma:2017pxb,Liu:2019urm}.
Many lattice works have been done on nucleon and meson PDFs, and generalized parton distributions (GPDs) based on the quasi-PDF approach~\cite{Lin:2013yra,Lin:2014zya,Chen:2016utp,Lin:2017ani,Alexandrou:2015rja,Alexandrou:2016jqi,Alexandrou:2017huk,Chen:2017mzz,Alexandrou:2018pbm,Chen:2018xof,Chen:2018fwa,Alexandrou:2018eet,Lin:2018qky,Fan:2018dxu,Liu:2018hxv,Wang:2019tgg,Lin:2019ocg,Chen:2019lcm,Lin:2020reh,Chai:2020nxw,Bhattacharya:2020cen,Lin:2020ssv,Zhang:2020dkn,Li:2020xml,Fan:2020nzz,Gao:2020ito,Lin:2020fsj,Zhang:2020rsx,Alexandrou:2020qtt,Alexandrou:2020zbe,Lin:2020rxa,Gao:2021hxl,Lin:2020rut}.
Alternative approaches to lightcone PDFs in lattice QCD are ``good lattice cross sections''~\cite{Ma:2017pxb,Bali:2017gfr,Bali:2018spj,Sufian:2019bol,Sufian:2020vzb} and the pseudo-PDF approach~\cite{Orginos:2017kos,Karpie:2017bzm,Karpie:2018zaz,Karpie:2019eiq,Joo:2019jct,Joo:2019bzr,Radyushkin:2018cvn,Zhang:2018ggy,Izubuchi:2018srq,Joo:2020spy,Bhat:2020ktg,Fan:2020cpa,Sufian:2020wcv,Karthik:2021qwz}.
However, LQCD calculations of the $x$ dependence of the pion PDFs have only been done for the valence-quark distribution~\cite{Chen:2018fwa,Sufian:2019bol,Izubuchi:2019lyk,Joo:2019bzr,Sufian:2020vzb,Shugert:2020tgq,Gao:2020ito}.

Only recently have lattice calculations of the gluon PDF become possible, when the necessary one-loop matching relations of the gluon PDF were computed for the pseudo-PDF~\cite{Balitsky:2019krf} and quasi-PDF~\cite{Zhang:2018diq,Wang:2019tgg} approaches.
Both approaches make direct calculation of the $x$ dependence of the pion gluon PDF feasible.
In this work, we apply the pseudo-PDF method by using the ratio renormalization scheme to avoid the difficulty of calculating the gluon renormalization factors.
There is a developed procedure for using the pseudo-PDF method to obtain lightcone PDFs from Ioffe-time distributions (ITDs) by matching through two steps, evolution and scheme conversion~\cite{Radyushkin:2018cvn}.
Using the pseudo-PDF method also allows us to use lattice correlators at all boost momenta at small Ioffe-time.
There have been a number of successful pseudo-PDF calculations of nucleon isovector
PDFs~\cite{Orginos:2017kos,Joo:2019jct,Joo:2020spy,Bhat:2020ktg} and pion valence-quark PDFs~\cite{Joo:2019bzr}.
The earliest calculation was done on a quenched lattice~\cite{Orginos:2017kos}, then the pion masses were set closer to the physical pion mass~\cite{Joo:2019jct,Joo:2019bzr,Joo:2020spy}, and the calculation at physical pion mass was done recently~\cite{Bhat:2020ktg}.
The lattice-calculated PDFs in Refs.~\cite{Joo:2019jct,Joo:2019bzr,Joo:2020spy,Bhat:2020ktg} show good agreement with the global-analysis PDFs.

In this work, we present the first calculation of the full $x$-dependent pion gluon distribution using the pseudo-PDF method from two lattice spacings, 0.12 and 0.15~fm, and three pion masses: 690, 310 and 220~MeV.
The rest of the paper is organized as follows.
In Sec.~\ref{sec:cal-details}, we present the procedure to obtain the lightcone gluon PDF from the reduced ITDs,
the numerical setup of lattice simulation, and how we extracted the reduced ITDs from lattice calculated correlators.
In Sec.~\ref{sec:results}, the final determination of the pion gluon PDF from our lattice calculations is compared with the NLO xFitter~\cite{Novikov:2020snp} and JAM pion gluon PDFs~\cite{Barry:2018ort,Cao:2021aci}.
The systematics induced by different steps are studied, and the lattice-spacing and pion-mass dependence are investigated.

